\subsection{The Hash and Equality Contract}

Hash tables route probes from \texttt{hash(key)} but cannot infer equality from hash equality alone. CPython therefore splits lookup into a cheap integer filter and an expensive object comparison.

\begin{itemize}
    \item \texttt{\_\_hash\_\_()}: produces the 64-bit probe seed stored in the table entry;
    \item \texttt{\_\_eq\_\_()}: resolves collisions and confirms identity of content.
\end{itemize}

The language invariant: if \texttt{a == b}, then \texttt{hash(a) == hash(b)}. The converse is false---collisions are expected and handled by probing.

\subsubsection{Key Uniqueness and Hash Invariants: The Coordination of \texttt{\_\_hash\_\_} and \texttt{\_\_eq\_\_}}

Dictionary and set lookup in CPython executes a strict two-stage gate at each probed slot:

\begin{enumerate}
    \item Compare the stored 64-bit hash field against \texttt{hash(search\_key)}. On mismatch, continue probing immediately---\texttt{\_\_eq\_\_} is not invoked.
    \item On hash match, dispatch \texttt{PyObject\_RichCompare} / \texttt{\_\_eq\_\_} between the stored key and the search key. Only a \texttt{True} result completes the lookup.
\end{enumerate}

This ordering keeps hot paths fast: most failed probes exit on integer comparison alone. Colliding but unequal objects pay the full rich-comparison cost.

The contract breaks catastrophically when a stored key mutates its hash-relevant state after insertion. The entry remains at the slot chosen from the \emph{old} hash; subsequent lookups using the \emph{new} hash probe a different region. The object becomes a permanent \emph{ghost entry}: still present in the table, unreachable by ordinary \texttt{in}, \texttt{dict[key]}, or \texttt{del}. Mutable built-ins (\texttt{list}, \texttt{dict}, \texttt{set}) set \texttt{\_\_hash\_\_ = None} to forbid this failure mode at construction time.

\subsubsection{The Hash Collision Vulnerability and Hash Randomization}

Deterministic string hashing would let an attacker craft distinct keys sharing the same bucket chain. Every insertion and lookup on that chain degrades to $O(n)$ work---a Hash-Flood denial-of-service against the interpreter event loop.

CPython mitigates this at process start by drawing a secret 64-bit seed and feeding it into the SipHash-2-4 family for \texttt{str} and \texttt{bytes} hashing. Identical string literals therefore produce different hash integers across separate interpreter processes, while remaining stable within one run. The table still uses open addressing locally; the global seed removes cross-process predictability of collision patterns.

Hash stability for user keys remains mandatory: any object acting as a set element or dict key must keep equality-relevant state fixed for the duration of storage.
